\renewcommand{\Anno}{2011}

\renewcommand{\NumeroProblema}{2}

\renewcommand{\Testo}{
\begin{enumerate}
	\item  Dimostrare che esistono due costanti reali $A$ e $B$ tali che
	$$\cos^{6}x+\sin^{6}x=A+B\cos(4x)$$
	per ogni numero reale $x$.

	\item  Siano $a$ ed $x$ numeri reali tali che $\cos^{6}x+\sin^{6}x=a$.
	
	Determinare, in funzione di $a$, il valore di $\cos^{8}x+\sin^{8}x$.

	\item  Determinare tutti i numeri reali $x\in[0,2\pi]$ tali che
	$$\cos^{6}x+\sin^{6}x=\frac{13}{16},$$ 
	precisando se si tratta di multipli razionali di $\pi$ (cioè se si scrivono nella forma $m\pi/n$ per opportuni interi $m$ ed $n$).
\end{enumerate}
}

\renewcommand{\Risposta}{
\begin{enumerate}
	\item[(a)]  Le costanti sono $A=5/8$ e $B=3/8$.

	\item[(b)]  $\displaystyle{\frac{2a^{2}+8a-1}{9}}$

	\item[(c)]  Nell'intervallo $[0,2\pi]$ l'equazione ha le seguenti 8 soluzioni
	$$\frac{\pi}{12},\quad\quad
	\frac{5\pi}{12},\quad\quad
	\frac{7\pi}{12},\quad\quad
	\frac{11\pi}{12},\quad\quad
	\frac{13\pi}{12},\quad\quad
	\frac{17\pi}{12},\quad\quad
	\frac{19\pi}{12},\quad\quad
	\frac{23\pi}{12},$$
	che sono tutte multipli razionali di $\pi$.
	
\end{enumerate}
}

\renewcommand{\Aiutino}{
\begin{enumerate}
	\renewcommand{\labelenumi}{(\alph{enumi})}
	\item[(a)]  Vi sono almeno 2 approcci possibili.
	\begin{itemize}
		\item Utilizzando la relazione trigonometrica fondamentale $\cos^{2}x+\sin^{2}x=1$, si può scrivere $\cos^{6}x+\sin^{6}x$ come somma di sole potenze di $\cos x$. Analogamente, utilizzando le opportune formule di duplicazione del coseno, si può scrivere $\cos(4x)$ come somma di sole potenze di $\cos x$. Non resta ora che imporre la relazione con coefficienti $A$ e $B$ incogniti e poi uguagliare i termini.
	
		\item Partendo dalla relazione 
		$$\cos^{6}x+\sin^{6}x = (\cos^{2}x+\sin^{2}x)^{3}-
		3\cos^{2}x\cdot\sin^{2}x\cdot(\cos^{2}x+\sin^{2}x),$$
		ci si riduce a $\cos^{2}x\cdot\sin^{2}x$, che è legato a $\sin^{2}(2x)$, il quale a sua volta è legato a $\cos(4x)$ da un'opportuna formula di duplicazione.
	\end{itemize}

	\item[(b)]  Partire dalla relazione
	$$\cos^{8}x+\sin^{8}x =  
	(\cos^{6}x+\sin^{6}x)(\cos^{2}x+\sin^{2}x)-\ldots$$
	
	Successivamente occorre procurarsi il valore di $\cos^{2}x\cdot\sin^{2}x$ in funzione di $a$, ma questo si può fare come illustrato nel punto precedente.

	\item[(c)] Ragionando come nei casi precedenti è possibile ricavarsi $\cos(4x)$ oppure $\cos^{2}x\cdot\sin^{2}x$.  Si arriva
	così ad una delle seguenti equazioni $$\cos(4x)=\frac{1}{2},
	\hspace{4em}
	\sin(2x)=\pm\frac{1}{2},$$
	da cui si conclude abbastanza agevolmente.  Nel trovare le soluzioni contenute in $[0,2\pi]$, si consiglia la massima attenzione alla gestione del termine ``$+2k\pi$'' nel momento in cui si va a dividere per 4 o per 2.
\end{enumerate}

}

\renewcommand{\SoluzioneUfficiale}{\TitoloSoluzione \DomandaA
Si ha che
\begin{eqnarray*}
	\cos^{6}x+\sin^{6}x & = & (\cos^{2}x+\sin^{2}x)^{3}-
	3\cos^{2}x\cdot\sin^{2}x\cdot(\cos^{2}x+\sin^{2}x)  \\
	 & = & 1-\frac{3}{4}\cdot 4\cos^{2}x\cdot\sin^{2}x  \\
	  & = & 1-\frac{3}{4}\sin^{2}(2x)  \\
	  & = & 1-\frac{3}{4}\cdot\frac{1-\cos(4x)}{2}  \\
	 & = & \frac{5}{8}+\frac{3}{8}\cos(4x),
\end{eqnarray*}
dove la penultima uguaglianza segue dalla relazione
$$\sin^{2}\alpha=\frac{1-\cos(2\alpha)}{2},$$
la quale a sua volta è conseguenza immediata della formula di duplicazione del coseno.

\DomandaB
Ragionando come nel punto precedente si ha che
\begin{eqnarray*}
	a & = & \cos^{6}x+\sin^{6}x  \\
	 & = & (\cos^{2}x+\sin^{2}x)^{3}-
	3\cos^{2}x\cdot\sin^{2}x\cdot(\cos^{2}x+\sin^{2}x)  \\
	 & = & 1-3\cos^{2}x\cdot\sin^{2}x,
\end{eqnarray*}
da cui
$$\cos^{2}x\cdot\sin^{2}x=\frac{1-a}{3}.$$

Pertanto si ha che
\begin{eqnarray*}
	\cos^{8}x+\sin^{8}x & = & (\cos^{6}x+\sin^{6}x)(\cos^{2}x+\sin^{2}x)
	-\cos^{2}x\cdot\sin^{2}x\cdot(\cos^{4}x+\sin^{4}x) \\
	 & = & \cos^{6}x+\sin^{6}x-\cos^{2}x\cdot\sin^{2}x\cdot\left[
	 (\cos^{2}x+\sin^{2}x)^{2}-2\cos^{2}x\cdot\sin^{2}x \right]  \\
	 & = & \cos^{6}x+\sin^{6}x-\cos^{2}x\cdot\sin^{2}x+2
	 \cos^{4}x\cdot\sin^{4}x\\
	 & = & a-\frac{1-a}{3}+2\frac{(1-a)^{2}}{9}  \\
	 & = & \frac{2a^{2}+8a-1}{9}.
\end{eqnarray*}

\DomandaC
Applicando la formula trovata al punto~(a) abbiamo che
$$\frac{5}{8}+\frac{3}{8}\cos(4x)=\frac{13}{16},$$
da cui si ricava facilmente che
$$\cos(4x)=\frac{1}{2}.$$

Abbiamo pertanto le 2 famiglie di soluzioni
$$4x=\pm\frac{\pi}{3}+2k\pi
\quad\quad
(k\in\mathbb{Z}),$$
da cui
$$x=\pm\frac{\pi}{12}+k\frac{\pi}{2}
\quad\quad
(k\in\mathbb{Z}).$$

Limitandoci all'intervallo $[0,2\pi]$ abbiamo pertanto le seguenti 8 soluzioni
$$\frac{\pi}{12},\quad\quad
\frac{5\pi}{12},\quad\quad
\frac{7\pi}{12},\quad\quad
\frac{11\pi}{12},\quad\quad
\frac{13\pi}{12},\quad\quad
\frac{17\pi}{12},\quad\quad
\frac{19\pi}{12},\quad\quad
\frac{23\pi}{12},$$
che sono chiaramente tutte multipli razionali di $\pi$.


}

\renewcommand{\SoluzioniAlternative}{\TitoloSoluzioneAlternativaA \DomandaA
Dalla relazione trigonometrica fondamentale abbiamo che
$$\cos^{6}x+\sin^{6}x=\cos^{6}x+(1-\cos^{2}x)^{3}=
1-3\cos^{2}x+3\cos^{4}x.$$

Applicando 2 volte la formula di duplicazione del coseno abbiamo invece che
$$\cos(4x)=2\cos^{2}(2x)-1=2(2\cos^{2}x-1)^{2}-1=8\cos^{4}x-8\cos^{2}x+1.$$

A questo punto la relazione
$$\cos^{6}x+\sin^{6}x=A+B\cos(4x)$$
equivale a
$$1-3\cos^{2}x+3\cos^{4}x=A+B-8B\cos^{2}x+8B\cos^{4}x.$$

Non resta ora che uguagliare i coefficienti di $\cos^{4}x$, $\cos^{2}x$, ed il termine noto.

\DomandaB
Consideriamo le 3 uguaglianze
\begin{eqnarray*}
	(\cos^{2}x+\sin^{2}x)^{2} & = & 
	\cos^{4}x+\sin^{4}x+2\cos^{2}x\cdot\sin^{2}x,  \\
	(\cos^{2}x+\sin^{2}x)^{3} & = &
	\cos^{6}x+\sin^{6}x+3\cos^{2}x\cdot\sin^{2}x\cdot
	(\cos^{2}x+\sin^{2}x),\\
	(\cos^{2}x+\sin^{2}x)^{4} & = &
	\cos^{8}x+\sin^{8}x+4\cos^{2}x\cdot\sin^{2}x\cdot
	(\cos^{4}x+\sin^{4}x)+6\cos^{4}x\cdot\sin^{4}x.
\end{eqnarray*}

In tutte le uguaglianze possiamo sostituire la relazione $\cos^{2}x+\sin^{2}x=1$.

In questo modo dalla seconda uguaglianza ricaviamo $\cos^{2}x\cdot\sin^{2}x$ in funzione di $a$.  A questo punto dalla prima possiamo ricavare $\cos^{4}x+\sin^{4}x$ in funzione di $a$. Ora, sfruttando la terza uguaglianza, abbiamo tutti gli elementi per ricavare $\cos^{8}x+\sin^{8}x$ in funzione di $a$.

\DomandaC
Ragionando come nel punto~(b) con $a=13/16$, si ottiene facilmente che
$$\cos^{2}x\cdot\sin^{2}x=\frac{1-a}{3}=\frac{1}{16}.$$

Per la formula di duplicazione del seno si ha quindi che
$$\sin^{2}(2x)=4\cos^{2}x\cdot\sin^{2}x=\frac{1}{4},$$
da cui
$$\sin(2x)=\pm\frac{1}{2}.$$

Abbiamo pertanto le seguenti 4 famiglie di soluzioni (si intende sempre al variare di $k\in\mathbb{Z}$)
$$2x=\pm\frac{\pi}{6}+2k\pi,
\hspace{5em}
2x=\pm\frac{5\pi}{6}+2k\pi,$$
da cui
$$x=\pm\frac{\pi}{12}+k\pi,
\hspace{5em}
x=\pm\frac{5\pi}{12}+k\pi.$$

Limitandoci all'intervallo $[0,2\pi]$ troviamo ovviamente le solite 8 soluzioni.
}

\renewcommand{\Osservazioni}{
}

\renewcommand{\Commenti}{
}